\chapter{Sviluppo pratico di un exploit}\label{chapter:exploit}
Finora abbiamo sempre avuto il codice sorgente a disposizione. Nel mondo reale però, molti servizi con cui è possibile interagire non mettono a disposizione né il codice sorgente né il file binario compilato. L'attaccante si trova quindi in uno scenario di tipo \emph{black-box}: non sa come sia scritto il programma, può solo fornirgli input e studiarne l'output o il comportamento.

Se però l'attaccante entra in possesso anche solo del binario, è possibile, tramite gli strumenti elencati nella Sottosezione~\ref{subsec:tools}, ottenere informazioni rilevanti tra cui: architettura, formato del file, impostazioni di sicurezza, disassemblato e decompilato (pseudocodice simile al linguaggio originale).

Tali informazioni rendono possibile ricostruire la logica del programma e individuarne le vulnerabilità: è esattamente ciò che tratta questo capitolo, insieme alla scrittura dello script Python che le sfrutta, applicando le tecniche del Capitolo~\ref{chapter:fmt_vuln}.


\section{Analisi del binario vulnerabile}\label{sec:vuln_service}
È necessario ottenere informazioni sul binario, di cui non sappiamo praticamente niente.
Il comando \texttt{file} permette di ricavare il formato del binario e la sua architettura (Codice~\ref{lst:file_output}) che sono rispettivamente ELF e x86\_64 poiché il binario è stato compilato con GCC nell'ambiente di riferimento in Tabella~\ref{tab:ambiente_riferimento}.
\begin{lstlisting}[
    language={},
    numbers=none,
    gobble=4,
    caption={Estratto dell'output del comando \texttt{file} sul binario},
    label={lst:file_output}
]
    ELF 64-bit LSB pie executable, x86-64, version 1 (SYSV), dynamically linked, [...], not stripped
\end{lstlisting}

Tipicamente i binari vengono compilati con l'opzione \texttt{-s} (\emph{strip}) che rimuove tutte le tabelle dei simboli e le informazioni di debug dal file eseguibile finale: oltre a ridurne le dimensioni, questo lo protegge rendendolo poco comprensibile. Per semplicità però, il binario in questo capitolo è stato compilato senza tale opzione. Vediamo infatti ``\texttt{not stripped}'' (Codice~\ref{lst:file_output}).

Un'altra informazione essenziale sono le impostazioni di sicurezza del binario che vengono ricavate dal comando \texttt{pwn checksec} (Codice~\ref{lst:checksec_output}).
\begin{lstlisting}[
    language={},
    numbers=none,
    gobble=4,
    caption={Output del comando \texttt{pwn checksec} sul binario},
    label={lst:checksec_output}
]
    Arch:       amd64-64-little
    RELRO:      Full RELRO
    Stack:      Canary found
    NX:         NX enabled
    PIE:        PIE enabled
    SHSTK:      Enabled
    IBT:        Enabled
    Stripped:   No
\end{lstlisting}
Tale configurazione è il default del compilatore GCC nel nostro ambiente (Tabella~\ref{tab:ambiente_riferimento}).

I comandi \texttt{file} e \texttt{pwn checksec} danno una prima idea del binario --- il tipo di disassemblato che si otterrà, le convenzioni in uso, le protezioni da affrontare --- ma nulla sulla logica del codice: per quella servono il disassemblato e il decompilato, entrambi ottenibili con \texttt{Ghidra}.

Il codice decompilato è un tentativo di ricostruire la logica e la struttura del programma originale esaminando il codice assembly (il disassemblato). Tale ricostruzione però è spesso confusionaria e va sistemata per rendere il codice più comprensibile (Codice~\ref{lst:ghidra_decompiled}).
\begin{lstlisting}[
    language=C,
    numbers=none,
    gobble=4,
    caption={Estratto del decompilato di \texttt{Ghidra} sistemato},
    label={lst:ghidra_decompiled}
]
    void debug_shell(void) {
        system("/bin/sh");
        return;
    }

    void log_message(char *msg) {
        printf(msg);
        return;
    }

    int main(void) {
        char buf[64];

        // ...
        puts("=== MiniLog v1.0 ===");
        while (fgets(buf, 64, stdin) != NULL) {
            log_message(buf);
        }
        // ...
        return 0;
    }
\end{lstlisting}
Il programma è una sorta di MiniLog che va a stampare in \texttt{stdout} tutto quello che viene scritto su \texttt{stdin}. Il problema sono le due funzioni \texttt{log\_message} e \texttt{debug\_shell}.

All'interno di \texttt{log\_message} è presente \texttt{printf(msg)} con \texttt{msg} controllata dall'utente e quindi vulnerabile.

La funzione \texttt{debug\_shell} apre una shell e, di per sé, non sarebbe un problema, poiché non viene mai chiamata da nessun punto del programma. In questo caso però, è pericolosa per la presenza della vulnerabilità format string che può dirottare il normale flusso d'esecuzione.

L'idea è di sfruttare la vulnerabilità format string andando a sovrascrivere il \emph{return address} (indirizzo di ritorno) di \texttt{log\_message} con l'indirizzo di \texttt{debug\_shell}: quando \texttt{log\_message} termina, l'esecuzione riprende da lì anziché da \texttt{main}, aprendo così una shell.
Per far ciò, è necessario conoscere sia l'indirizzo della funzione \texttt{debug\_shell} sia l'indirizzo dello stack dove è salvato il return address di \texttt{log\_message}.

Questi indirizzi però vengono randomizzati a causa di PIE (Codice~\ref{lst:checksec_output}) e ASLR, quest'ultima attiva di default a livello di sistema operativo, che complicano ma non rendono impossibile lo sfruttamento, soprattutto grazie alla presenza del ciclo \texttt{while}, che ci permette di iterare più volte sulla \texttt{printf} vulnerabile.

È infatti possibile calcolare l'indirizzo di \texttt{debug\_shell} a partire dal leak di un indirizzo dell'eseguibile, e calcolare l'indirizzo dove è salvato il return address di \texttt{log\_message} a partire dal leak di un indirizzo dello stack. In una successiva iterazione è poi possibile sovrascrivere tale return address con l'indirizzo di \texttt{debug\_shell} appena calcolato.


\section{Fase di leak: individuazione degli indirizzi utili}\label{sec:leak_phase}

